\section{The true adjoint torsion density}\label{sec:torsion}
We use positive cochain torsion \(T\), with the convention that an invertible \(0\)-to-\(1\) differential contributes the absolute determinant. Its acyclic analytic convention is \(\prod_q\det'(\Delta_q)^{(-1)^{q+1}q/2}\). For a closed irreducible flat local system in dimension three, take the tangent basis in \(H^1\) and its dual in \(H^2\) under Poincar\'e duality. Then \(T^{1/2}\) is a density in tangent coordinates: changing the \(H^1\) basis and its \(H^2\) dual contributes the square of the tangent determinant before taking the square root. We use the torsion gluing rule for a short exact sequence of based complexes, including the torsion of the long cohomology sequence \cite{Milnor}. The closed torsion has the equivalent Ray--Singer norm by Cheeger--M\"uller \cite{CM1,CM2}; no analytic gluing formula for pieces with boundary is needed.

\begin{theorem}[Determinant-line identification on every cover]\label{thm:torsion}
On the regular branch \(h_A=-1\), \(h_B\) noncentral, the pushforward of the actual positive adjoint torsion density \(T^{1/2}\) to the exterior radii \(x_B\), summing the regular bending and cap-root sheets, is \(C_{G,A}\,dM_B\). Here \(C_{G,A}>0\) depends on normalization and the cut, not on \(x_B\). The statement holds for a minimum cover with \(A\)--\(A\) seams as well as for an independent cover. The regular branch has dimension \(3N-2\nu-3\).
\end{theorem}

\begin{proof}
We compute the cochain determinants on the pieces, the long cohomology sequence on the cut tori, and finally the Poincar\'e pairing that turns the result into a tangent density.

\paragraph{\(B\) block, valence \(k\).}
Because its fibre \(h_b\) is noncentral, the adjoint system preserves an axial line and its perpendicular plane. All the block holonomies fix the axis; in particular the axial local system is trivial. Remove both exceptional solid tori. The remaining bundle is \(\Sigma_{0,k+2}\times S^1\); the plane circle differential \(\Ad(h_b)-I\) is invertible. Filter the product cochain complex by cellular base degree. Each quotient is the circle complex with that differential, with its torsion sign shifted by its cell degree. Their product gives plane torsion \(|\det(\Ad(h_b)-I)|^{\chi(\Sigma)}=|\det(\Ad(h_b)-I)|^{-k}\). The cap core is the section generator \(U_b\) or \(V_b\), up to inversion, in the actual primitive \((2,\pm1)\) filling slopes. Its solid torus contributes \(|\det(\Ad(U_b)-I)|\) or \(|\det(\Ad(V_b)-I)|\) to the plane torsion. Because the sections of \(B\) are central, \(V_b=\pm U_b^{-1}\). The plane on each boundary torus has identical nonzero determinants in successive degrees, hence torsion one and no cohomology. Here all determinants are restricted to the plane. The multiplicative rule therefore gives exactly
\begin{equation}\label{eq:B}
 T_{\rm plane}(B_b)^{1/2}
 =\frac{|\det(\Ad(U_b)-I)|}{|\det(\Ad(h_b)-I)|^{k/2}}
 =\frac{4\cos^2(\pi x_b/2)}{(2\sin(\pi x_b))^k}.
\end{equation}
The axial cohomology and determinants are constant in the marked slopes. Its \(H^0\) is \(\RR\). Its \(H^1\) restricts at the \(k\) tori to a common fibre value and independently a sum-zero section vector; the opposite exceptional slopes cancel in the base relation. Its \(H^2\) restricts to a sum-zero vector of oriented area classes. These descriptions also follow directly from ordinary cohomology: in real coefficients the two fillings express the exceptional section periods as opposite halves of the fibre period, and the outer boundary cycles generate the planar base homology subject to their sum.

\paragraph{\(A\) block.}
Its adjoint fibre holonomy is identity and its base representation is irreducible. After removing the exceptional solid tori, write \(W'=H^1(\Sigma;\Ad)\), where \(\Sigma\) includes all boundary circles. \(H^0(\Sigma;\Ad)=0\) and \(H^2(\Sigma;\Ad)=0\). The product has \(H^1=W'\), \(H^2=\eta\wedge W'\), where \(\eta\) evaluates as \(1\) on a regular fibre. In paired bases the two shifted copies of the surface complex have cancelling torsions. The boundary trace map from \(W'\) to the sum of the boundary invariant Lie algebras is onto by the relative sequence: \(H^2(\Sigma,\partial\Sigma;\Ad)=H^0(\Sigma;\Ad)^*=0\). These invariant algebras are axial lines at the caps and the \(A\)--\(B\) sections, and \(\RR^3\) at \(A\)--\(A\) sections.

For each removed cap, its core has adjoint rotation \(\pi\), so all nonzero plane torsion factors are constant. In the axial cohomology sequence the \(H^0\) of the solid torus maps isomorphically to that of its boundary. Its \(H^1\) maps with nonzero fibre coefficient of magnitude two in the marked slope bases, while the surface deformation maps to the independent section direction. The resulting degree-one differential onto the seam cohomology is onto. The degree-two restriction from the unfilled piece maps \(\eta\wedge W'\) onto the cap trace lines. These cohomology determinants are fixed slope determinants in splitting bases for the surjective boundary trace map. Thus the filled \(A\) piece has \(H^1=W_a\), \(H^2=\eta\wedge W_a\), \(H^0=H^3=0\), where \(W_a\subset W'\) has both cap traces fixed. The positive torsion of the piece is constant in paired bases. The real class \(\eta\) may be adjusted to extend across the caps by splitting the ordinary first cohomology with fibre value \(1\); the base has boundary. Its horizontal adjustment does not change the stated \(H^2\) description because absolute \(H^2(\Sigma;\Ad)=0\).

Write \(W_a^0\) for the kernel of the remaining seam restriction. These are the fixed-class tangent modes of the punctured sphere with its two cap classes and its noncentral boundary classes fixed; a central boundary class is a point and its deformation restricts by zero in the full \(\RR^3\). If \(k_a\) is the number of \(B\) neighbours, \(\dim W_a^0=2k_a-2\): there are \(k_a+2\) two-dimensional conjugacy classes, and the regular product constraint and effective conjugacy quotient each subtract three. Let \(\Omega_a\) be the Goldman symplectic form on that fixed-class space.

\paragraph{Seams and the closed long exact sequence.}
At an \(A\)--\(B\) seam only the axial cohomology survives, with dimensions \(1,2,1\). The plane is acyclic with torsion one. In oriented coordinates \((A\text{ fibre},A\text{ section})=(B\text{ section},B\text{ fibre})\), the \(A\) degree-one image is in the section direction, and the \(B\) image has sum-zero \(A\)-fibre coordinates incident to that \(B\), together with a common \(A\)-section value. At an \(A\)--\(A\) seam both \(\SU\) holonomies are central. Its adjoint torus complex is ordinary cohomology tensored with \(\RR^3\). The two \(A\) degree-one section maps go to complementary three-dimensional summands under the swap. They are onto their summands simultaneously with the other boundary maps by the surjectivity above. Matching therefore forces both section restrictions to vanish and yields no new degree-one cokernel at that seam. The \(H^0(T_e)=\RR^3\) survives in the gluing tangent space, and the difference of the two \(A\) degree-two restrictions is onto \(H^2(T_e)=\RR^3\); its kernel keeps one common \(\RR^3\) datum.

In more detail, let \(E_{AB},E_{AA}\) be the two edge sets. The global \(H^0,H^3\) vanish. The connecting image of the seam \(H^0\)'s in \(H^1(M_G;\Ad)\) is
\[
 L=\bigl(\RR^{E_{AB}}/\text{one diagonal per }b\in B\bigr)
       \ \oplus\ \bigoplus_{e\in E_{AA}}\RR^3 .
\]
The remaining \(H^1\) is the sum of the \(B\) radii and \(\oplus_a W_a^0\): matching forces all the \(B\) section vectors to be zero and the \(A\) section traces at \(A\)--\(B\) seams to be their \(B\) fibre periods. The connecting image of the seam \(H^1\)'s in \(H^2(M_G;\Ad)\) is one axial quotient per \(b\), given by the \(A\)-fibre vectors modulo the \(B\) sum-zero vectors. The remaining \(H^2\) consists of \(\oplus_a W_a^0\) in their \(\eta\)-copies, one \(B\) sum-zero area vector at every \(B\), and the common degree-two \(\RR^3\) at every \(A\)--\(A\) seam. To obtain these descriptions choose the trace splitting bases at \(A\) using the onto restriction map. The whole cohomology correction has constant matrices in these bases. Rotating between orthonormal adjoint frames to use the actual seam identifications changes their absolute determinants by one. In particular the number of tangent directions is
\[
 |B|+(|E_{AB}|-|B|)+3|E_{AA}|+\sum_{a\in A}(2k_a-2)
  =3N-2\nu-3.
\]

\paragraph{Dualizing the determinant line.}
The identity pairing a Mayer--Vietoris connecting class by restriction and cup product on the cut torus gives a triangular description of the closed Poincar\'e pairing. The axial degree-two connecting quotient pairs with the \(B\) fibre radii (their derivatives in the unit axis are \(\pm\pi\,dx_b\)); it is zero on the fixed-class and bending modes. The bending modes represented by \(L\) pair with the sum-zero axial degree-two data and with the common \(\RR^3\)'s, respectively, by constant quotient pairings. For the remaining pairing, use relative representatives of \(W_a^0\) with zero boundary trace and extend them by zero from that block. Pairing with \(\eta\) times the second representative integrates \(\eta\) to one along the regular fibre, leaving the fixed-class relative surface cup product \(\Omega_a\). The representatives can be chosen away from the cap neighbourhoods; no term from a cap or another block changes this middle pairing. Unspecified off-diagonal pairings do not change the triangular determinant. Replacing the \(H^2\) bases with Poincar\'e-dual ones therefore multiplies the positive torsion by \(\prod_a|\det\Omega_a|\) up to a constant. The square root supplies \(\prod_a|\operatorname{pf}\Omega_a|\), the surface Liouville measures.

The bending modes are actual gluing parameters. At each \(A\)--\(B\) seam their centralizer is \(U(1)\), with a common \(U(1)\) quotient at each \(B\); at an \(A\)--\(A\) seam it is \(\SU\) up to the finite centres. The tangent bases used above are unit-axis or orthonormal adjoint bases for these groups. Their volumes and finite quotients are independent of the radii in the regular chambers. A tree has no further loop condition on these identifications: fixing the piece representations at the prescribed classes leaves the indicated choices of edge intertwiners, modulo the vertex stabilizers. The two cap-root sheets per \(B\) also have fixed multiplicity.

\paragraph{Cancellation of the seam sine factors.}
For a fixed-class \(\SU\) conjugacy class of angle \(\theta\), the quasi-Hamiltonian Liouville measure is proportional to \(|\det(\Ad g-I)_{\perp}|^{1/2}\) times its normalized Haar measure, hence proportional to \(\sin\theta\) for noncentral classes, with a common constant \cite[Prop.~3.6]{AMW}. A central class is a point. The pushforward under fusion is convolution, and its value at the regular product level is proportional to the reduced symplectic volume \cite[\S5.1]{AMW}. These results apply at the regular radii specified above, since all configurations at the product level are irreducible there. The two cap sine factors are constants. Thus the volume at \(a\) is proportional to
\[
 \prod_{b\sim a,\ b\in B}\sin(\pi x_b)\ L_a(x).
\]
One can equally obtain this factor from the usual character formula for punctured-sphere volumes; at degree one the sum is understood in the heat or Abel sense. A central product shift and inversions are immaterial since \(K(-g)=K(g^{-1})=K(g)\). Each cross seam supplies exactly one sine factor. Together they cancel every sine denominator in \eqref{eq:B}. The remaining density is a constant times \(\prod_b2\cos^2(\pi x_b/2)\,dx_b\prod_aL_a(x)\), which is \eqref{eq:haar}. This proves the assertion, including the central--central case.
\end{proof}

\paragraph{Two-site check.}
The smallest path has two Klein-bottle interval bundles sewn by the fibre--section swap. On the irreducible branch with the left fibre central and the right one regular, the adjoint system splits globally into the right axis and its plane. In the Klein presentation \(XYX^{-1}=Y^{-1}\), the plane Fox differentials are \(d_0(v)=((X-I)v,(Y-I)v)\) and \(d_1(a,b)=(I-Y^{-1})a+(X+Y^{-1})b\). The left plane and the seam torus have torsion one by cancellation of the determinants with noncentral \(Y\). On the right \(Y=I\), so the plane square-root torsion is \(\sqrt{|\det(X-I)|/|\det(X+I)|}=\cot(\pi x/2)\). The axial matrices and the \(dx\) duality normalization are constant. This directly checks both endpoint behaviours in \eqref{eq:haar}, which here gives \(\frac{\pi}{4}\cot(\pi x/2)\,dx\).
